\documentclass[11pt]{article}
\usepackage[T1]{fontenc}
\usepackage{lmodern,amsmath,amssymb,amsthm,mathtools,microtype}
\usepackage[margin=1in]{geometry}
\usepackage{enumitem,booktabs,longtable,array}
\usepackage[colorlinks=true,linkcolor=blue,citecolor=blue,urlcolor=blue]{hyperref}
\setlist[enumerate]{leftmargin=*,itemsep=4pt}
\setlength{\emergencystretch}{2em}
\newtheorem{theorem}{Theorem}[section]
\newtheorem{proposition}[theorem]{Proposition}
\newtheorem{lemma}[theorem]{Lemma}
\newtheorem{corollary}[theorem]{Corollary}
\theoremstyle{definition}
\newtheorem{definition}[theorem]{Definition}
\theoremstyle{remark}
\newtheorem{remark}[theorem]{Remark}
\newcommand{\C}{\mathbb C}
\newcommand{\R}{\mathbb R}
\newcommand{\Pp}{\mathbb P}
\newcommand{\Dscr}{\mathcal D'}
\newcommand{\M}{\mathcal M}
\newcommand{\len}{\operatorname{len}}
\newcommand{\crit}{\mathrm{crit}}
\newcommand{\Dom}{\mathsf D}
\newcommand{\dd}{\,\mathrm d}
\newcommand{\code}[1]{{\small\nolinkurl{#1}}}
\newcommand{\leanref}[2]{\href{https://github.com/leobon12/lean-formalizations/blob/de92cf632a47ee74bf4aa407c9e5e5b9b71c3ded/SupercriticalConformalCovariance/#1}{#2}}
\title{Conformal covariance of the supercritical\\Liouville quantum gravity metric}
\author{}
\date{October 25, 2026}
\begin{document}
\maketitle
\begin{abstract}
We give a domain-field formulation of conformal covariance for the
supercritical Liouville quantum gravity metric. For a zero-boundary Gaussian
free field plus a possibly dependent continuous function on a connected
Greenian domain, the metric has a measurable representative depending only
on the field. Under a fixed conformal change of coordinates, the metrics
agree almost surely at every pair of points, including singular points.
Representatives can be chosen for any prescribed countable family of field
laws and maps. We organize the result in the form of Gwynne--Miller's
conformal-covariance theorem while retaining the quantifiers of the
accompanying Lean statement. The proof derives this formulation from
realization-based covariance and realization independence, with an explicit
construction of domain metrics and descent to the original probability space.
\end{abstract}

\section{Introduction}
Fix $\xi>\xi_{\crit}$ and let $Q=Q(\xi)\in(0,2)$ be the LFPP distance
exponent. Let $H\mapsto D_H$ be a strong whole-plane LQG metric with
parameter $\xi$ and one fixed normalization. Here \emph{strong} refers
to Ding--Gwynne's axioms \cite[Definition 1.7]{DG}: the length-space
property, locality, continuous Weyl scaling, and translation and dilation
covariance. Their existence and uniqueness theorem identifies this
assignment up to a deterministic positive factor. We use the same
assignment throughout.

Distances may take the value $+\infty$. A domain means a nonempty
connected open subset of $\C$. The main results concern Greenian
domains, on which the zero-boundary GFF is defined, without any
boundary-regularity assumption. Write $\Dscr(U)$ for distributions on
$U$ with the measurable structure generated by test-function evaluations.
Set $\M(U)=[0,\infty]^{U\times U}$ with its product measurable structure.
The representatives below are measurable maps into $\M(U)$; their
values are required to be metrics almost surely under the specified
field laws, not at every distribution.

\begin{definition}[Domain fields]\label{def:fields}
An admissible field on a Greenian domain $U$ is a measurable random
distribution $h$ on $(\Omega,\mathcal F,\Pp)$ with
\[
                  h=h_0^U+f\quad\text{almost surely},
\]
where $h_0^U$ is a zero-boundary GFF on $U$ and $f$ is a measurable
random element of $C(U,\R)$. The decomposition is supplied on the
given probability space. The function $f$ may depend on $h_0^U$,
need not be bounded, and need not extend continuously to $\partial U$.
\end{definition}

A whole-plane realization of $h$ is a pair $(H,F)$ with $H$ a
whole-plane GFF plus a continuous function and $F$ a measurable
random element of $C(U,\R)$ such that $h=H|_U+F$ almost surely.
Realizations are also allowed on measure-preserving extensions
$\pi:\Omega'\to\Omega$, with $h\circ\pi=H|_U+F$.
For a realization define
\begin{equation}\label{eq:domain}
 \Dom^U_{H,F}(x,y)=
 \inf_{P:x\to y,\;P\subset U}
       \int_0^{\len(P;D_H)}e^{\xi F(P(t))}\dd t .
\end{equation}
Paths are rectifiable and parametrized by $D_H$-length.
An empty infimum is $+\infty$, and a constant path has length zero.
This is an \emph{internal} distance in $U$.

\begin{definition}[A version for a specified field]\label{def:version}
A measurable $\Phi_U:\Dscr(U)\to\M(U)$ is a metric version for $h$
if, on every measure-preserving extension carrying a realization
$(H,F)$ of $h\circ\pi$,
\begin{equation}\label{eq:version}
                  \Phi_U(h\circ\pi)=\Dom^U_{H,F}
             \quad\text{almost surely as functions on }U^2.
\end{equation}
Whenever such a version is asserted, existence of a realization is
included. The null event may depend on the extension and realization.
\end{definition}

For a deterministic conformal bijection $\phi:U\to V$, write
\begin{equation}\label{eq:transform}
            h^\phi=h\circ\phi^{-1}+Q\log|(\phi^{-1})'|.
\end{equation}
Composition of distributions uses pullback of test functions with the
appropriate Jacobian.

\begin{theorem}[Conformal covariance]\label{thm:main}
Let $U$ be a connected Greenian domain, let $\phi:U\to V$ be a fixed
deterministic conformal bijection, and let $h$ be an admissible field
on $U$. There exist measurable maps
\[
       \Phi_U:\Dscr(U)\to\M(U),\qquad
       \Phi_V:\Dscr(V)\to\M(V)
\]
which are metric versions for $h$ and $h^\phi$, respectively.
Both fields have whole-plane realizations on product extensions of
their original probability space. Almost surely on
$(\Omega,\mathcal F,\Pp)$,
\begin{equation}\label{eq:main}
 \Phi_U(h)(x,y)=\Phi_V(h^\phi)(\phi(x),\phi(y))
                         \qquad\text{for all }x,y\in U.
\end{equation}
Thus, with $D_h^U=\Phi_U(h)$ and $D_{h^\phi}^V=\Phi_V(h^\phi)$,
\[
                  D_h^U(x,y)=D_{h^\phi}^V(\phi(x),\phi(y)).
\]
\end{theorem}

This is the forward convention of Gwynne--Miller's Axiom IV and
Theorem 1.3 \cite{GM}. The version and existence clauses identify the
metrics with the fixed whole-plane construction and exclude vacuous
choices of representatives. Equality is simultaneous in the endpoints,
including singular endpoints, with exact constant one. The event may
depend on the fixed map, parameter, and field law.

\begin{theorem}[A prescribed countable family]\label{thm:countable}
Keep $\xi,Q,D,U,V$ fixed. For each index $i$ in a countable set $I$,
let $\phi_i:U\to V$ be a deterministic conformal bijection and let
$h_i$ be admissible on $(\Omega_i,\mathcal F_i,\Pp_i)$.
One pair $\Phi_U,\Phi_V$ can be chosen before the index $i$ so that,
for every $i$, both product realizations exist, the maps are versions
for $h_i$ and $h_i^{\phi_i}$ on every realization extension, and
\[
 \Pp_i\!\left[
 \Phi_U(h_i)(x,y)=
 \Phi_V(h_i^{\phi_i})(\phi_i(x),\phi_i(y))
                  \text{ for all }x,y\in U\right]=1.
\]
\end{theorem}

The representatives are chosen after the law, or prescribed countable
family of laws, is specified. We do not assert one $\Phi_U$ chosen
before \emph{all} admissible laws, or a common sample event for an
uncountable collection of conformal maps.

\subsection*{Proof organization and sources}
We construct the domain metric on a realization, represent it measurably
in terms of the field, and then transfer realization-based covariance to
the original probability space. A positive-weight mixture supplies common
representatives for a countable family.

The analytic covariance input is the released Lean theorem
\cite{Lean}, stated precisely in Proposition~\ref{prop:root}.
This rewrite proves the domain-field formulation from that theorem;
it does not reproduce its entire analytic comparison proof as an
independent proof. Section~\ref{sec:analytic} explains that proof's
structure and gives a repaired endpoint estimate.
Appendix~\ref{app:lean} records the declaration correspondence.

Gwynne--Miller's theorem supplies the form of the statement, not the
supercritical input. Devlin's simultaneous coordinate-change theorem
\cite[Theorem 1.1]{Devlin} assumes $\xi<\xi_{\crit}$ and is not used here.

\section{Domain metrics and realizations}\label{sec:domain}

\begin{lemma}\label{lem:realize}
An admissible field on $U$ has a whole-plane realization on a
measure-preserving product extension
$\Omega\times\Dscr(\C)\to\Omega$. The transformed field $h^\phi$ is
admissible on $V$ and has such a realization as well.
\end{lemma}
\begin{proof}
The whole-plane Markov decomposition on $U$ gives a reference coupling
\[
                     H|_U=Z+\mathfrak h_U,
\]
where $Z$ has the zero-boundary GFF law and $\mathfrak h_U$ has a
continuous harmonic version on $U$. Couple the reference whole-plane
field to the given $h_0^U$ using its conditional law given $Z$.
This keeps the full joint law of $(h,h_0^U,f)$ unchanged. On the
product extension,
\[
               h\circ\pi=H|_U+(f\circ\pi-\mathfrak h_U).
\]
The correction is continuous. No independence between $f$ and $h_0^U$
is required.

Conformal invariance of the zero-boundary GFF gives
\[
 h^\phi=(h_0^U\circ\phi^{-1})
          +(f\circ\phi^{-1}+Q\log|(\phi^{-1})'|).
\]
The first summand is a zero-boundary GFF on $V$ and the second is
a measurable random continuous function. The image $V$ is connected
and Greenian. Apply the same realization construction on $V$.
\end{proof}

The corresponding product construction is
\leanref{GMForm/Base.lean}{\code{exists_prod_realization_of_zeroBoundary}}.
It includes the requisite measurable continuous versions.

\subsection{The domain forms of the metric axioms}

Formula \eqref{eq:domain} is the intrinsic length metric obtained by
weighting $D_H$-length. For every open $W\subset U$ it satisfies
\begin{equation}\label{eq:restriction}
       \Dom^W_{H,F|_W}(x,y)=\Dom^U_{H,F}(x,y;W).
\end{equation}
For $g\in C(U,\R)$ it satisfies
\begin{equation}\label{eq:weyl}
                     \Dom^U_{H,F+g}
                         =e^{\xi g}\cdot\Dom^U_{H,F},
\end{equation}
where the right-hand side is path-length reweighting internal to $U$.
These identities hold on the realization-level sample event
simultaneously in continuous corrections.

To justify functions unbounded near $\partial U$, observe that every
path used in a distance has compact image inside $U$. Continuous
weights are bounded above and below on that image. The weighted
length formula applies there, successive reweighting multiplies the
weights, and infimizing over paths gives
\eqref{eq:restriction}--\eqref{eq:weyl}.
If $F$ extends to $\widehat F\in C(\C)$, then
\[
                 \Dom^U_{H,F}=D_{H+\widehat F}(\cdot,\cdot;U).
\]
Otherwise this identity is applied only on relatively compact
subsets using cutoff extensions.

\begin{proposition}[Realization independence]\label{prop:ind}
If $(H,F)$ and $(\widehat H,\widehat F)$ realize the same domain
field on one probability space, then
\[
                  \Dom^U_{H,F}=\Dom^U_{\widehat H,\widehat F}
                        \quad\text{almost surely on }U^2.
\]
This statement applies on any open set for which the realizations
are supplied.
\end{proposition}

This is the released theorem
\leanref{Frozen/DomainMetric_ae_eq_of_isDomainRealization.lean}
{\code{domainMetric_ae_eq_of_isDomainRealization}}.
It is a pathwise identification on the coupling, not merely equality
in distribution. Its proof uses locality for cutoff whole-plane
fields on a countable compact exhaustion, followed by the internal
path construction. In particular, fields agreeing on an open overlap
have equal internal metrics on that overlap. The domain length,
restriction, locality, and Weyl statements are developed in
\leanref{Domain/Axioms.lean}{\code{Domain/Axioms.lean}}.
They are the domain counterparts of the first three axioms in
Gwynne--Miller's presentation. They do not assert a single measurable
representative valid before every field law is specified.

\section{The analytic covariance theorem}\label{sec:analytic}

\begin{proposition}[Covariance for realized fields]\label{prop:root}
Let $\xi>\xi_{\crit}$, $Q=Q(\xi)$, and $D$ be any strong whole-plane
LQG metric with these parameters. Let $U,V$ be connected Greenian
domains and $\psi:U\to V$ a fixed deterministic conformal bijection.
If $k$ on $V$ has a realization $(H,F)$, the field
$k\circ\psi+Q\log|\psi'|$ has a realization on a
measure-preserving extension. On every such extension and for every
such realization $(G,K)$, almost surely,
\begin{equation}\label{eq:root}
 \Dom^V_{H\circ\pi,F\circ\pi}(\psi(u),\psi(v))
                         =\Dom^U_{G,K}(u,v)
                         \quad\text{for all }u,v\in U.
\end{equation}
\end{proposition}

We use the frozen
\leanref{Frozen/ConformalCovariance.lean}{\code{conformal_covariance}}
theorem of \cite{Lean}. The existence clause is part of the theorem.
It imposes no additional connector or comparison assumption on a
strong metric.

Here is the structure of its analytic proof. First compare the source
metric $d$ with the conformally pulled-back metric $e$.
Annular estimates and a filled-set loop chain give a deterministic
finite bi-Lipschitz constant. The filled-set construction ensures
connectivity and includes the terminal loop's cost. Small conformal
maps are then compared with affine tangents, using a smooth field
normalization, compact confinement, and lower semicontinuous metric
convergence. For any sequence of fields, a countable mixture supplies
one locality functional valid for all its laws. The confinement event
uses a strict annular barrier margin.

Recovery in the lower semicontinuous topology moves endpoints.
The following deterministic estimate explains how short loops recover
the actual endpoints, with the radial cost fully accounted for.

\newpage
\begin{lemma}[Two-scale endpoint recovery]\label{lem:recovery}
Suppose $d_n,e_n$ are extended length metrics on domains eventually
containing every compact subset of $\C$. Their finite-length paths
are Euclidean continuous, $C^{-1}d_n\le e_n\le Cd_n$, and both
sequences converge locally in the lower semicontinuous topology to
the same extended metric $d$.

Let $u_n\to u$, $v_n\to v$ with $u\ne v$, and suppose
$L_n=d_n(u_n,v_n)\to L<\infty$. Suppose $s_j\downarrow0$ and,
for $x=u_n,v_n$, every fixed $j$, and all sufficiently large $n$,
\begin{equation}\label{eq:regularity}
 d_n(x,\partial B_{s_j}(x))\le Ms_j^\theta,\qquad
 d_n(\mathrm{around}\,A_{s_j,\sqrt{s_j}}(x))\le Ms_j^\theta
\end{equation}
for $0<M<\infty$ and $\theta>0$. Then
\[
                   \limsup_n e_n(u_n,v_n)\le d(u,v)\le L.
\]
\end{lemma}
\begin{proof}
The lower-bound clause of convergence gives $d(u,v)\le L$.
Choose recovery endpoints $a_n\to u$, $b_n\to v$ for $e_n$
and paths $Q_n$ between them whose $e_n$-lengths tend to $d(u,v)$.
Fix $k$ large enough that the $s_k$-balls about the two limiting
endpoints are disjoint. Choose $j$ with $\sqrt{s_j}<s_k$.

For large $n$, the recovery endpoints lie inside the corresponding
$s_j$-balls. Use the first bound in \eqref{eq:regularity} at scale
$s_k$ to reach $\partial B_{s_k}(u_n)$ at $d_n$-cost at most
$2Ms_k^\theta$. Use the second bound at scale $s_j$ to find a
surrounding loop of length at most $2Ms_j^\theta$.
The radial path crosses the loop since $s_k>\sqrt{s_j}$. Following
the radial path to its first intersection and then a portion of the
loop connects $u_n$ to any prescribed point of the loop at cost
at most $2M(s_k^\theta+s_j^\theta)$. Repeat at $v_n$.

Pointwise metric comparison compares path lengths by taking suprema
over partitions. The path $Q_n$ crosses both loops by Euclidean
continuity. Splicing its middle part to the connectors gives
\[
 \limsup_n e_n(u_n,v_n)
       \le d(u,v)+4CM(s_k^\theta+s_j^\theta).
\]
Choose $j=j(k)\to\infty$ as $k\to\infty$ and let $k\to\infty$.
\end{proof}

A short path only to radius $s_j$ need not reach a loop in
$A_{s_j,\sqrt{s_j}}$. The larger-radius estimate is therefore
essential. This repairs the connector step of the earlier manuscript.

The analytic argument next builds good annuli whose regular geodesic
crossings have comparison coefficient $1+\eta$.
Near-independence across annular scales and a finite lattice cover
give an event valid for all relevant geodesics. A definite fraction
of each retained interval lies in a good crossing. If $C_*$ is the
optimal upper comparison constant, this yields
\[
 C_{\mathrm{new}}=(1+\eta)+\frac{A}{A+1}(C_*-1-\eta)<C_*,
                \qquad 0<\eta<C_*-1 .
\]
The retained intervals are selected using a last-exit/first-entry
corridor so that the geodesic endpoints lie outside the enlarged
balls. The contradiction gives $C_*=1$; the inverse-map argument
gives equality. The complete analytic proof, including the probability
estimates, is in the \code{BiLip}, \code{Tangent}, \code{Recovery},
\code{Cover}, \code{Improve}, and \code{Root} modules of \cite{Lean}.
This summary is not a replacement for those proofs.

\section{Measurable metric versions}\label{sec:versions}

\begin{lemma}[A representative for one realization]\label{lem:functional}
For a specified field $h$ with a realization $(H,F)$ on $U$, there
is a measurable $\Phi_U:\Dscr(U)\to\M(U)$ such that
$\Phi_U(h)=\Dom^U_{H,F}$ almost surely as functions on $U^2$.
\end{lemma}
\begin{proof}
Choose an increasing relatively compact open exhaustion $W_n\Subset U$
which eventually contains every compact subset of $U$. Choose a
measurable cutoff extension $F_n\in C(\C)$ equal to $F$ on $W_n$:
multiply by a deterministic continuous cutoff compactly supported
in $U$, and extend by zero.

Locality of the whole-plane metric gives a measurable
$A_n:\Dscr(W_n)\to\M(W_n)$ representing
$D_{H+F_n}(\cdot,\cdot;W_n)$ as a function of
$(H+F_n)|_{W_n}=h|_{W_n}$. This representative is chosen for the
law under consideration. Define
\begin{equation}\label{eq:measurable}
 \Phi_U(k)(x,y)=
       \inf_{\substack{n\ge1\\x,y\in W_n}}
                    A_n(k|_{W_n})(x,y).
\end{equation}
For each pair this is a countable infimum of measurable functions,
so the map into the evaluation-product space $\M(U)$ is measurable.

Intersect the countably many locality and Weyl events. On this event,
each term is the internal weighted distance on $W_n$, simultaneously
for all its endpoint pairs. The infimum equals \eqref{eq:domain},
since every path in $U$ has compact image and is contained in some
$W_n$. The resulting equality holds as a function on all of $U^2$.
\end{proof}

This is
\leanref{Domain/Axioms.lean}{\code{exists_measurable_domainMetric}}.
In particular, no unjustified factorization into an uncountable
product is needed.

\begin{lemma}[Descent through a product extension]\label{lem:descent}
Let $\mu$ be a probability measure on $\Omega\times X$ with first
marginal $\Pp$, where $X$ is a nonempty standard Borel space.
If a predicate $\mathcal A(\omega)$ holds for $\mu$-almost every
$(\omega,x)$, then it holds for $\Pp$-almost every $\omega$,
with almost-everywhere statements understood in the completed measure.
\end{lemma}
\begin{proof}
Choose a measurable $\mu$-null set $N$ containing all pairs where
$\mathcal A(\omega)$ fails. Disintegrate $\mu$ over its first
marginal with conditional probabilities $K(\omega,\dd x)$.
For almost every $\omega$, $K(\omega,N_\omega)=0$.
If $\mathcal A(\omega)$ fails, $N_\omega=X$ and this probability
equals one. Thus the failing set lies in a measurable $\Pp$-null set.
\end{proof}

This is \leanref{GMForm/Base.lean}{\code{ae_of_ae_fst}}.
It applies to the predicate asserting equality at every endpoint
pair, without first assuming measurability of that uncountable
intersection.

\begin{lemma}[Every realization extension]\label{lem:all}
If a measurable $\Phi_U$ agrees with the metric on one product
realization of $h$, then it agrees with the metric on every
realization on every measure-preserving extension of the same base field.
\end{lemma}
\begin{proof}
Disintegrate the fixed product realization over $\Omega$.
Given another realization on $\pi:\Omega'\to\Omega$, sample the
auxiliary whole-plane field by this conditional law at $\pi(\omega')$.
The resulting common refinement on $\Omega'\times\Dscr(\C)$
projects with the correct laws to both realizations.

On the refinement, pull back the equality defining $\Phi_U$.
Proposition~\ref{prop:ind} identifies its realized metric with the
competing realized metric. The desired equality therefore holds
almost surely on the refinement. It depends only on $\omega'$,
so Lemma~\ref{lem:descent} descends it to $\Omega'$.
\end{proof}

The original probability space need not be standard Borel.
The auxiliary distribution space supplies the conditional kernel.
These constructions are
\leanref{GMForm/Base.lean}{\code{product_common_refinement}} and
\leanref{GMForm/Base.lean}{\code{version_on_all_extensions}}.

\section{Proof of Theorem 1.3}

By Lemma~\ref{lem:realize}, choose product realizations of $h$ and
$h^\phi$ over their original probability space.
Lemma~\ref{lem:functional} supplies measurable representatives
$\Phi_U,\Phi_V$. Lemma~\ref{lem:all} makes them versions on every
realization extension, as required by Definition~\ref{def:version}.

Apply Proposition~\ref{prop:root} with the domains interchanged
and with $\phi^{-1}:V\to U$. The field on its target is $h$;
the pulled-back field on its source is
$h\circ\phi^{-1}+Q\log|(\phi^{-1})'|=h^\phi$.
On a further product extension carrying the target realization,
\[
             \Dom^U_{H,F}(x,y)
                  =\Dom^V_{G,K}(\phi(x),\phi(y))
                        \quad\text{for all }x,y\in U
\]
almost surely. All fields are understood to be pulled to the common
extension. The version clauses replace these two realized metrics
by $\Phi_U(h)$ and $\Phi_V(h^\phi)$.

The resulting all-endpoints identity depends only on the original
sample $\omega$. Apply Lemma~\ref{lem:descent} through the two product
extensions to obtain \eqref{eq:main} on the original probability
space. This proves Theorem~\ref{thm:main}.

Every equality used here takes values in $[0,\infty]$; no infinite
distance is subtracted or divided. Singular endpoints are therefore
retained. The whole-plane assignment and its normalization are fixed
on both sides, so there is no extra multiplicative constant.

\section{Countably many laws and maps}\label{sec:countable}

\begin{lemma}\label{lem:mixture}
For any prescribed countable family of realized fields
$h_i=H_i|_U+F_i$, one measurable map $\Phi_U$ represents all their
domain metrics almost surely under their respective laws.
\end{lemma}
\begin{proof}
For a nonempty family choose positive weights $w_i$ summing to one
and form the disjoint-union law $\sum_iw_i\Pp_i$. On component $i$
set $(h,H,F)=(h_i,H_i,F_i)$. Normalize the GFF component of each $H_i$
in the same way, absorbing any additive constant into its continuous
correction. The mixture GFF component has this same law, while the
continuous correction may depend on both the field and the index.
Thus the mixture is again a whole-plane realization.

Apply Lemma~\ref{lem:functional} once to this mixture. Every component
has positive weight, so its almost-sure identity holds under every
$\Pp_i$, with the same $\Phi_U$. An empty family imposes no condition.
\end{proof}

\begin{proof}[Proof of Theorem~\ref{thm:countable}]
Construct product realizations for $h_i$ and $h_i^{\phi_i}$ over each
$\Omega_i$. Apply Lemma~\ref{lem:mixture} separately to the source
and target families to choose one pair $\Phi_U,\Phi_V$.
For every $i$, Lemma~\ref{lem:all} makes these representatives
versions on every realization extension. The proof of
Theorem~\ref{thm:main}, with the fixed pair of representatives,
now applies to $h_i$ and $\phi_i$ and gives equality on $\Omega_i$.
\end{proof}

The mixture selects functions; it does not replace the original
probability spaces by a single asserted sample space. If a countable
family does share one probability space, its full-probability events
can be intersected. A countable-mixture construction does not establish
one representative for every admissible law.

\clearpage
\appendix
\section{Correspondence and scope}\label{app:lean}
This rewrite uses public revision
\href{https://github.com/leobon12/lean-formalizations/commit/de92cf632a47ee74bf4aa407c9e5e5b9b71c3ded}{\texttt{de92cf632a47}} of
\href{https://github.com/leobon12/lean-formalizations}{the accompanying Lean repository}.
Paths below are relative to \code{SupercriticalConformalCovariance}.

\begin{longtable}{@{}p{0.26\textwidth}p{0.68\textwidth}@{}}
\toprule
Manuscript & Released declaration or module\\ \midrule
\endhead
Theorem~\ref{thm:main} &
\leanref{GMForm/Base.lean}{\code{GMForm.gm_axiomIV_on_base}}\\[5pt]
Theorem~\ref{thm:countable} &
\leanref{GMForm/Base.lean}{\code{GMForm.exists_countably_universal_gm_axiomIV_on_base}}\\[5pt]
Proposition~\ref{prop:root} &
\leanref{Frozen/ConformalCovariance.lean}{\code{conformal_covariance}}\\[5pt]
Proposition~\ref{prop:ind} &
\leanref{Frozen/DomainMetric_ae_eq_of_isDomainRealization.lean}{\code{domainMetric_ae_eq_of_isDomainRealization}}\\[5pt]
Domain construction &
\leanref{Frozen/DomainMetric.lean}{\code{Frozen/DomainMetric.lean}};
\leanref{Domain/Axioms.lean}{\code{Domain/Axioms.lean}}\\[5pt]
Product realization &
\leanref{GMForm/Base.lean}{\code{GMForm.exists_prod_realization_of_zeroBoundary}}\\[5pt]
Versions on extensions &
\leanref{GMForm/Base.lean}{\code{GMForm.version_on_all_extensions}}\\[5pt]
Base-space descent &
\leanref{GMForm/Base.lean}{\code{GMForm.ae_of_ae_fst}}\\[5pt]
Countable representative &
\leanref{GMForm/Countable.lean}{\code{GMForm.exists_countably_universal_domain_metric}}\\
\bottomrule
\end{longtable}

The main theorems retain the Lean statements' connected Greenian
domains, supplied zero-boundary decomposition, measurable random
continuous correction, fixed parameter, fixed normalization, and
versions on every realization extension. Their equality is on the
field's original probability space and is simultaneous in all endpoints.

The frozen covariance and realization-independence theorems are the
analytic inputs. The domain-field theorems are supporting declarations
derived from them. The release guide distinguishes their audit status
from the frozen theorem seals. This prose rewrite is not a new kernel
build. The separate
\leanref{Root/AllXi.lean}{\code{Root.conformal_covariance_all_xi}}
theorem removes the supercritical restriction for the realization-based
statement. We keep $\xi>\xi_{\crit}$ in Theorems~\ref{thm:main} and
\ref{thm:countable}, matching their \code{GMForm} declarations.

The earlier manuscript is superseded by this domain-field formulation.
Its radial connector is repaired in Lemma~\ref{lem:recovery}; the other
analytic repairs are supplied by the released covariance theorem.
The law-specific and countable-family scope remains distinct from one
representative chosen before all admissible laws or all conformal maps.

\begin{thebibliography}{9}
\bibitem{DG}
J. Ding and E. Gwynne.
\emph{Uniqueness of the critical and supercritical Liouville quantum
gravity metrics}. Proc. London Math. Soc. 126 (2023), 216--333.
\url{https://arxiv.org/abs/2110.00177}.
\bibitem{GM}
E. Gwynne and J. Miller.
\emph{Conformal covariance of the Liouville quantum gravity metric
for $\gamma\in(0,2)$}. Ann. Inst. H. Poincar\'e Probab. Statist.
57 (2021), 1016--1031.
\url{https://arxiv.org/abs/1905.00384}.
\bibitem{Devlin}
C. Devlin VI.
\emph{The coordinate change formula for the Liouville quantum gravity
metric holds for all conformal maps simultaneously}. 2026.
\url{https://arxiv.org/abs/2603.04640}.
\bibitem{Lean}
\emph{Conformal covariance of the supercritical Liouville quantum
gravity metric}. Accompanying Lean development and statement guide,
public revision \texttt{de92cf632a47} (2026).
\href{https://github.com/leobon12/lean-formalizations/tree/de92cf632a47ee74bf4aa407c9e5e5b9b71c3ded/results/supercritical-conformal-covariance}{Source and statement guide at the cited revision}.
\end{thebibliography}
\end{document}
